%! AP Statistics — Unit 1, Lesson 1.2 — Slides (Beamer)
\documentclass[aspectratio=169, 11pt]{beamer}
\usepackage{apstats-beamer}
\usetikzlibrary{shapes.geometric, arrows.meta, positioning, fit, backgrounds}

%=============================================================================
\begin{document}

% ─────────────────────────────────────────────────────────────────────────────
% SLIDE 1 — LESSON TITLE
% ─────────────────────────────────────────────────────────────────────────────
{
\setbeamercolor{background canvas}{bg=navy}
\begin{frame}[plain]
  \begin{tikzpicture}[remember picture, overlay]
    \fill[navylight] (current page.north west)
      rectangle ([yshift=-1.35cm]current page.north east);
  \end{tikzpicture}
  \vspace{2.8cm}
  \hspace{0.55cm}%
  \begin{minipage}{0.9\textwidth}
    {\color{goldacc}\bfseries\small LESSON 1.2}\\[0.5cm]
    {\color{white}\bfseries\LARGE The Language of Variation:\\[0.25cm]
      Variables}\\[1.4cm]
    {\color{skymid}\small AP Statistics~$\cdot$~Unit 1: Exploring One-Variable Data}
  \end{minipage}
\end{frame}
}

% ─────────────────────────────────────────────────────────────────────────────
% SLIDE 2 — PRIMARY OBJECTIVE
% ─────────────────────────────────────────────────────────────────────────────
{
\setbeamercolor{background canvas}{bg=sky}
\begin{frame}[t, plain]
  \begin{tikzpicture}[remember picture, overlay]
    \fill[navy] (current page.north west)
      rectangle ([yshift=-0.25cm]current page.north east);
  \end{tikzpicture}
  \vspace{0.45cm}\par
  \hspace{0.55cm}%
  \begin{minipage}{0.9\textwidth}

    \sectionlabel{Primary Objective}

    \normalsize
    Students will classify variables as \textbf{categorical} or
    \textbf{quantitative}, and further distinguish quantitative variables as
    \textbf{discrete} or \textbf{continuous}. Students will articulate why
    variable type determines which statistical methods are appropriate
    (Big Idea: VAR).

    \vspace{0.55cm}

    \normalsize\textbf{\color{navy}By the end of this lesson, I will be able to:}

    \smallskip
    \begin{itemize}
      \setlength{\itemsep}{0.3em}\setlength{\topsep}{0.2em}
      \item Classify a variable as categorical or quantitative and justify
            using the ``average test.''
      \item Distinguish between discrete and continuous quantitative variables
            and give an example of each.
      \item Explain why variable type must be identified before choosing a
            statistical display or summary.
    \end{itemize}

  \end{minipage}
\end{frame}
}

% ─────────────────────────────────────────────────────────────────────────────
% SLIDE 3 — HOOK: SORT THESE VARIABLES
% ─────────────────────────────────────────────────────────────────────────────
{
\setbeamercolor{background canvas}{bg=hookbg}
\begin{frame}[t, plain]
  \begin{tikzpicture}[remember picture, overlay]
    \fill[goldacc] (current page.north west)
      rectangle ([yshift=-1.35cm]current page.north east);
    \node[anchor=west, text=white, font=\bfseries\large,
          inner xsep=0.55cm, inner ysep=0pt]
      at ([yshift=-0.68cm]current page.north west)
      {Hook --- Sort These Variables};
  \end{tikzpicture}
  \vspace{1.1cm}\par
  \hspace{0.55cm}%
  \begin{minipage}{0.9\textwidth}
    \small
    {\color{white}\bfseries
      The table below shows data recorded for four students.
      Sort the five variables into two groups --- without looking at any labels yet.
      What was your sorting rule?}

    \vspace{0.4cm}

    \setlength{\tabcolsep}{8pt}
    \renewcommand{\arraystretch}{1.5}
    \begin{tabular}{@{} >{\bfseries\color{white}}p{0.14\textwidth}
                        >{\centering\arraybackslash}p{0.10\textwidth}
                        >{\centering\arraybackslash}p{0.10\textwidth}
                        >{\centering\arraybackslash}p{0.14\textwidth}
                        >{\centering\arraybackslash}p{0.14\textwidth} @{}}
      \rowcolor{navy}
      {\color{skymid}Name} &
      {\color{skymid}Grade} &
      {\color{skymid}GPA} &
      {\color{skymid}Fav.\ Color} &
      {\color{skymid}Sleep (hrs)} \\
      \rowcolor{white!15!hookbg}
      {\color{white}Jordan}  & {\color{white}11} & {\color{white}3.7} &
        {\color{white}Blue}   & {\color{white}7.5} \\
      \rowcolor{white!8!hookbg}
      {\color{white}Priya}   & {\color{white}10} & {\color{white}3.2} &
        {\color{white}Green}  & {\color{white}6.0} \\
      \rowcolor{white!15!hookbg}
      {\color{white}Marcus}  & {\color{white}12} & {\color{white}2.9} &
        {\color{white}Red}    & {\color{white}8.0} \\
      \rowcolor{white!8!hookbg}
      {\color{white}Sofia}   & {\color{white}11} & {\color{white}3.5} &
        {\color{white}Purple} & {\color{white}7.0} \\
    \end{tabular}

    \vspace{0.45cm}
    {\color{white}\bfseries\small
      \textbf{Group A:}\;\underline{\hspace{5cm}} \qquad
      \textbf{Group B:}\;\underline{\hspace{5cm}}}
  \end{minipage}
\end{frame}
}

% ─────────────────────────────────────────────────────────────────────────────
% SLIDE 4 — VOCABULARY & KEY CONCEPTS
% ─────────────────────────────────────────────────────────────────────────────
{
\setbeamercolor{background canvas}{bg=greenbg}
\begin{frame}[t, plain]
  \begin{tikzpicture}[remember picture, overlay]
    \fill[greenacc] (current page.north west)
      rectangle ([yshift=-1.35cm]current page.north east);
    \node[anchor=west, text=white, font=\bfseries\large,
          inner xsep=0.55cm]
      at ([yshift=-0.68cm]current page.north west)
      {Vocabulary \& Key Concepts};
  \end{tikzpicture}
  \vspace{1.1cm}\par
  \hspace{0.4cm}%
  \begin{minipage}{0.93\textwidth}
    \vspace{0.4em}
    \footnotesize
    \begin{multicols}{2}
      \textbf{\color{greenacc}Variable} --- a characteristic that can take
        different values for different individuals\\[5pt]
      \textbf{\color{greenacc}Categorical variable} --- places an individual
        into one of several groups or categories; no meaningful arithmetic\\[5pt]
      \textbf{\color{greenacc}Quantitative variable} --- takes numerical values
        for which arithmetic makes sense

      \columnbreak

      \textbf{\color{greenacc}Discrete quantitative variable} --- countable
        values with gaps (e.g., number of siblings: 0, 1, 2, \ldots)\\[5pt]
      \textbf{\color{greenacc}Continuous quantitative variable} --- can take
        any value in an interval; measured, not counted (e.g., height, time)\\[5pt]
      \textbf{\color{greenacc}Distribution of a variable} --- describes what
        values a variable takes and how often; developed in Lessons 1.3--1.10
    \end{multicols}
  \end{minipage}
\end{frame}
}

% ─────────────────────────────────────────────────────────────────────────────
% SLIDE 5 — CATEGORICAL VS. QUANTITATIVE: THE AVERAGE TEST
% ─────────────────────────────────────────────────────────────────────────────
{
\setbeamercolor{background canvas}{bg=white}
\begin{frame}[t, plain]
  \navyheader{Categorical vs.\ Quantitative: The Average Test}
  \hspace{0.4cm}%
  \begin{minipage}{0.93\textwidth}
    \vspace{0.2em}
    \begin{columns}[T]

      \begin{column}{0.48\textwidth}
        \small
        \textbf{\color{greenacc}Categorical Variable}\\[4pt]
        Places an individual into a \textbf{group or category}.
        The values are words or labels --- or numbers that act as labels.\\[6pt]
        \textit{Examples:} favorite color, mode of transport,
        blood type, ZIP code\\[10pt]

        \textbf{\color{navy}Quantitative Variable}\\[4pt]
        Takes \textbf{numerical values} for which arithmetic makes sense.
        You can add, subtract, and average the values meaningfully.\\[6pt]
        \textit{Examples:} GPA, hours of sleep, commute time, salary
      \end{column}

      \begin{column}{0.48\textwidth}
        \small
        \textbf{\color{goldacc}The Average Test}\\[4pt]
        Ask: \textit{``Does the average of these values give a
        \textbf{meaningful} result in context?''}

        \vspace{8pt}
        \setlength{\tabcolsep}{6pt}
        \renewcommand{\arraystretch}{1.6}
        \begin{tabular}{@{} >{\bfseries}p{0.12\textwidth} p{0.72\textwidth} @{}}
          \rowcolor{greenbg}
          Yes & The variable is \textbf{quantitative}. \\
          \rowcolor{redbg}
          No  & The variable is \textbf{categorical}. \\
        \end{tabular}

        \vspace{10pt}
        \footnotesize\color{slate}
        \textit{``Does averaging jersey numbers tell you anything about
        the players?''}\\[4pt]
        \textit{``Does averaging commute times tell you something useful?''}
      \end{column}

    \end{columns}
  \end{minipage}
\end{frame}
}

% ─────────────────────────────────────────────────────────────────────────────
% SLIDE 6 — TRAP VARIABLES
% ─────────────────────────────────────────────────────────────────────────────
{
\setbeamercolor{background canvas}{bg=white}
\begin{frame}[t, plain]
  \navyheader{Watch Out: Numbers That Are Categorical}
  \hspace{0.4cm}%
  \begin{minipage}{0.93\textwidth}
    \vspace{0.3em}
    \small
    \textbf{\color{redacc}A number is not automatically quantitative.}
    If the number is being used as a \textit{label} or \textit{identifier},
    it is categorical --- averaging it gives a meaningless result.

    \vspace{0.6em}

    \setlength{\tabcolsep}{8pt}
    \renewcommand{\arraystretch}{1.7}
    \begin{tabular}{@{} >{\bfseries\color{navy}}p{0.22\textwidth}
                        p{0.36\textwidth}
                        p{0.32\textwidth} @{}}
      \rowcolor{sky}
      Variable & Why it looks quantitative & Why it is categorical \\
      \midrule
      ZIP code       & It's a 5-digit number  & Labels a geographic area; average ZIP = nonsense \\
      Jersey number  & It's a number on a uniform & Identifies a player; no meaningful arithmetic \\
      Student ID     & It's a numeric code    & A unique identifier; averaging IDs means nothing \\
      Phone area code & It's a 3-digit number & Labels a calling region; not a measured amount \\
      Likert rating (1--5) & Looks like a scale & Equal spacing not guaranteed; context-dependent \\
    \end{tabular}

    \vspace{0.5em}
    \footnotesize\color{slate}
    \textit{The same \textbf{question} can yield different variable types.
    ``How do students rate the cafeteria?''
    can be recorded as a 1--5 scale \textbf{or} as ``satisfied / unsatisfied'' ---
    your choice of measurement determines the variable type.}
  \end{minipage}
\end{frame}
}

% ─────────────────────────────────────────────────────────────────────────────
% SLIDE 7 — CLASSIFYING VARIABLES: EXAMPLES
% ─────────────────────────────────────────────────────────────────────────────
{
\setbeamercolor{background canvas}{bg=white}
\begin{frame}[t, plain]
  \navyheader{Classifying Variables: Examples}
  \hspace{0.4cm}%
  \begin{minipage}{0.93\textwidth}
    \vspace{0.3em}
    \small
    Classify each variable: \textbf{C} = Categorical,
    \textbf{QD} = Quantitative Discrete,
    \textbf{QC} = Quantitative Continuous.

    \vspace{0.5em}
    \setlength{\tabcolsep}{7pt}
    \renewcommand{\arraystretch}{1.65}
    \begin{tabular}{@{} c >{\raggedright\arraybackslash}p{0.38\textwidth}
                        >{\centering\arraybackslash}p{0.08\textwidth}
                        p{0.35\textwidth} @{}}
      \rowcolor{sky}
      \# & Variable & Type & Reason \\
      \midrule
      1 & Eye color                        & C   & Labels a category; no arithmetic \\
      2 & Number of siblings               & QD  & Countable; gaps between values \\
      3 & Body temperature (\textdegree F) & QC  & Measured; can be any decimal \\
      4 & Area code                        & C   & Numeric label; averaging meaningless \\
      5 & Number of words in an essay      & QD  & Count of whole words \\
      6 & Time to run a mile (min)         & QC  & Measured time; any positive value \\
      7 & Blood type (A, B, AB, O)         & C   & Places individual into a category \\
      8 & Annual salary (\$)               & QC  & Measured amount; treated as continuous \\
    \end{tabular}
  \end{minipage}
\end{frame}
}

% ─────────────────────────────────────────────────────────────────────────────
% SLIDE 8 — DISCRETE VS. CONTINUOUS
% ─────────────────────────────────────────────────────────────────────────────
{
\setbeamercolor{background canvas}{bg=white}
\begin{frame}[t, plain]
  \navyheader{Quantitative Variables: Discrete vs.\ Continuous}
  \hspace{0.4cm}%
  \begin{minipage}{0.93\textwidth}
    \vspace{0.3em}
    \begin{columns}[T]

      \begin{column}{0.47\textwidth}
        \small
        \textbf{\color{navy}Discrete}\\[4pt]
        Values are \textbf{countable} --- gaps exist between them.
        You can list all possible values.\\[6pt]
        \textit{Ask: ``Can I count them?''}\\[6pt]
        \begin{tabular}{@{} >{\color{greenacc}\bfseries}l l @{}}
          Anchor: & Number of pets \\
          Values: & 0, 1, 2, 3, \ldots \\
        \end{tabular}\\[6pt]
        \footnotesize\color{slate}
        More examples: number of siblings, songs on a playlist,
        students in a class
      \end{column}

      \begin{column}{0.47\textwidth}
        \small
        \textbf{\color{navy}Continuous}\\[4pt]
        Can take \textbf{any value in an interval} --- measured, not counted.
        No gaps; values can fall between any two others.\\[6pt]
        \textit{Ask: ``Could it fall between two values?''}\\[6pt]
        \begin{tabular}{@{} >{\color{redacc}\bfseries}l l @{}}
          Anchor: & Weight of a pet \\
          Values: & any positive real number \\
        \end{tabular}\\[6pt]
        \footnotesize\color{slate}
        More examples: height, commute time, temperature,
        song duration
      \end{column}

    \end{columns}

    \vspace{0.6em}
    \small
    \setlength{\tabcolsep}{7pt}
    \renewcommand{\arraystretch}{1.5}
    \begin{tabular}{@{} p{0.34\textwidth}
                        >{\centering\arraybackslash}p{0.14\textwidth}
                        >{\centering\arraybackslash}p{0.14\textwidth} @{}}
      \rowcolor{navy}
      {\color{white}\bfseries Question} &
      {\color{white}\bfseries Discrete} &
      {\color{white}\bfseries Continuous} \\
      Are there gaps between values? & Yes & No \\
      \rowcolor{sky}
      Is it measured on a scale?     & No  & Yes \\
      Can I list all possible values? & Yes & No \\
    \end{tabular}
  \end{minipage}
\end{frame}
}

% ─────────────────────────────────────────────────────────────────────────────
% SLIDE 9 — DECISION FLOWCHART
% ─────────────────────────────────────────────────────────────────────────────
{
\setbeamercolor{background canvas}{bg=white}
\begin{frame}[plain]
  \navyheader{Decision Flowchart: Classifying Any Variable}

  \vspace{-0.3cm}
  \begin{center}
  \begin{tikzpicture}[
    node distance = 0.55cm and 1.4cm,
    decision/.style = {
      diamond, draw=navy, fill=sky, thick,
      text width=3.0cm, align=center,
      inner sep=2pt, aspect=1.8
    },
    terminal/.style = {
      rectangle, rounded corners=4pt,
      draw=navy, thick,
      text width=2.4cm, align=center,
      minimum height=0.8cm, font=\bfseries\small
    },
    cat/.style    = {terminal, fill=greenbg,  draw=greenacc, text=greenacc},
    discr/.style  = {terminal, fill=redbg,    draw=redacc,   text=redacc},
    cont/.style   = {terminal, fill=goldbg,   draw=goldacc,  text=goldacc},
    arrow/.style  = {-Stealth, thick, draw=navy},
    lbl/.style    = {font=\footnotesize\bfseries, text=navy}
  ]

  % Nodes
  \node[decision] (q1)
    {Is the variable a word, label, or category?};

  \node[cat, below left=0.7cm and 2.2cm of q1] (cat1)
    {Categorical};

  \node[decision, below=1.1cm of q1] (q2)
    {Does averaging give a \emph{meaningful} result?};

  \node[cat, below left=0.7cm and 2.2cm of q2] (cat2)
    {Categorical \\ \footnotesize\normalfont (trap number)};

  \node[decision, below=1.1cm of q2] (q3)
    {Can you \emph{count} the values? (Are there gaps?)};

  \node[discr, below left=0.7cm and 2.2cm of q3] (qd)
    {Quantitative \\ Discrete};

  \node[cont, below right=0.7cm and 2.2cm of q3] (qc)
    {Quantitative \\ Continuous};

  % Arrows
  \draw[arrow] (q1) -- node[lbl, above left=-2pt]{Yes} (cat1);
  \draw[arrow] (q1) -- node[lbl, right=2pt]{No}  (q2);
  \draw[arrow] (q2) -- node[lbl, above left=-2pt]{No}  (cat2);
  \draw[arrow] (q2) -- node[lbl, right=2pt]{Yes} (q3);
  \draw[arrow] (q3) -- node[lbl, above left=-2pt]{Yes} (qd);
  \draw[arrow] (q3) -- node[lbl, above right=-2pt]{No}  (qc);

  \end{tikzpicture}
  \end{center}
\end{frame}
}

% ─────────────────────────────────────────────────────────────────────────────
% SLIDE 10 — ACTIVE MONITORING: MINI-WHITEBOARD ROUND
% ─────────────────────────────────────────────────────────────────────────────
{
\setbeamercolor{background canvas}{bg=sky}
\begin{frame}[t, plain]
  \begin{tikzpicture}[remember picture, overlay]
    \fill[navy] (current page.north west)
      rectangle ([yshift=-1.35cm]current page.north east);
    \node[anchor=west, text=white, font=\bfseries\large,
          inner xsep=0.55cm, inner ysep=0pt]
      at ([yshift=-0.68cm]current page.north west)
      {Active Monitoring --- Mini-Whiteboard Round};
  \end{tikzpicture}
  \vspace{1.1cm}\par
  \hspace{0.55cm}%
  \begin{minipage}{0.9\textwidth}
    \vspace{0.4em}
    \small

    For each variable your teacher displays, hold up your whiteboard showing:\\[6pt]
    \hspace{1cm}
    \textbf{\color{greenacc}C} \hspace{0.8cm}
    \textbf{\color{navy}QD} \hspace{0.8cm}
    \textbf{\color{redacc}QC}\\[8pt]
    Then give a one-phrase reason out loud.

    \vspace{0.6cm}

    \begin{tabular}{@{} >{\bfseries\color{navy}}p{0.12\textwidth}
                        p{0.80\textwidth} @{}}
      Tip 1: & Use the average test first --- does averaging give a meaningful result? \\[4pt]
      Tip 2: & If quantitative, ask ``Can I count it?'' to decide discrete vs.\ continuous. \\[4pt]
      Tip 3: & Watch for numbers that are really labels. \\
    \end{tabular}

    \vspace{0.7cm}
    \textcolor{slate}{\small\itshape
      See Part 1 of the activity worksheet for the practice problems.}
  \end{minipage}
\end{frame}
}

% ─────────────────────────────────────────────────────────────────────────────
% SLIDE 11 — GROUP ACTIVITY: CARD SORT
% ─────────────────────────────────────────────────────────────────────────────
{
\setbeamercolor{background canvas}{bg=redbg}
\begin{frame}[t, plain]
  \begin{tikzpicture}[remember picture, overlay]
    \fill[redacc] (current page.north west)
      rectangle ([yshift=-1.35cm]current page.north east);
    \node[anchor=west, text=white, font=\bfseries\large,
          inner xsep=0.55cm, inner ysep=0pt]
      at ([yshift=-0.68cm]current page.north west)
      {Group Activity --- Variable Card Sort};
  \end{tikzpicture}
  \vspace{1.1cm}\par
  \hspace{0.55cm}%
  \begin{minipage}{0.9\textwidth}
    \vspace{0.4em}
    \small

    \begin{enumerate}
      \setlength{\itemsep}{0.5em}\setlength{\topsep}{0em}
      \item Open your activity worksheet to \textbf{Part 2}.
      \item Sort the variable cards into three columns on the worksheet:\\[4pt]
        \hspace{0.8cm}
        \textbf{\color{greenacc}Categorical} \hspace{1.2cm}
        \textbf{\color{navy}Quantitative Discrete} \hspace{1.2cm}
        \textbf{\color{redacc}Quantitative Continuous}
      \item When your group finishes, compare your sort with a neighboring group.
            Resolve any disagreements before debriefing.
      \item Be ready to defend the cards your group found most contested.
    \end{enumerate}

    \vspace{0.7cm}
    {\color{redacc}\bfseries\small TIME: 8--10 minutes}

    \vspace{0.3cm}
    \textcolor{slate}{\small\itshape
      Your teacher will debrief 2--3 contested cards whole-class.}
  \end{minipage}
\end{frame}
}

% ─────────────────────────────────────────────────────────────────────────────
% SLIDE 12 — EXIT TICKET
% ─────────────────────────────────────────────────────────────────────────────
{
\setbeamercolor{background canvas}{bg=navy}
\begin{frame}[t, plain]
  \begin{tikzpicture}[remember picture, overlay]
    \fill[navylight] (current page.north west)
      rectangle ([yshift=-1.35cm]current page.north east);
  \end{tikzpicture}
  \vspace{0.45cm}\par
  \hspace{0.55cm}%
  \begin{minipage}{0.9\textwidth}
    {\color{goldacc}\bfseries\footnotesize\MakeUppercase{Exit Ticket}}
    \vspace{0.5cm}\par
    \small

    {\color{skymid} A researcher is studying music listening habits.
    For each variable on your exit ticket, write the type and a
    one-phrase justification.}

    \vspace{0.5cm}
    \begin{enumerate}
      \setlength{\itemsep}{0.5em}\setlength{\topsep}{0em}
      \item {\color{white}Classify each of the five variables on your exit ticket
            as \textbf{C}, \textbf{QD}, or \textbf{QC}.
            Write a brief justification for each.}
      \item {\color{white}Answer the AP-style multiple-choice question on
            annual income and job type.
            Write your answer \textit{and} explain your reasoning in one sentence.}
    \end{enumerate}

    \vspace{0.6cm}
    {\color{skymid}\small\itshape
      Completeness is graded. Justification is required for full credit.}
  \end{minipage}
\end{frame}
}

% ─────────────────────────────────────────────────────────────────────────────
% SLIDE 13 — CLOSING & HOMEWORK
% ─────────────────────────────────────────────────────────────────────────────
{
\setbeamercolor{background canvas}{bg=navy}
\begin{frame}[t, plain]
  \begin{tikzpicture}[remember picture, overlay]
    \fill[navylight] (current page.north west)
      rectangle ([yshift=-1.35cm]current page.north east);
  \end{tikzpicture}
  \vspace{0.45cm}\par
  \hspace{0.55cm}%
  \begin{minipage}{0.9\textwidth}
    {\color{goldacc}\bfseries\footnotesize\MakeUppercase{Closing \& Homework}}
    \vspace{0.4cm}\par
    \begin{columns}[T]

      \begin{column}{0.54\textwidth}
        \small
        {\color{goldacc}\bfseries Key Reminders}\\[5pt]
        \begin{itemize}
          \setlength{\itemsep}{0.4em}\setlength{\topsep}{0em}
          \item {\color{white}\textbf{Numbers are not always quantitative.}
            {\color{skymid} ZIP codes, jersey numbers, and IDs are
            categorical labels.}}
          \item {\color{white}\textbf{Use the average test.}
            {\color{skymid} ``Does averaging these values give a
            meaningful result?''}}
          \item {\color{white}\textbf{Discrete = countable; continuous = measured.}
            {\color{skymid} Ask: ``Could it fall between two values?''}}
          \item {\color{white}\textbf{Variable type is Step 0.}
            {\color{skymid} Identify it before choosing any display
            or summary statistic.}}
        \end{itemize}
      \end{column}

      \begin{column}{0.42\textwidth}
        \small
        {\color{goldacc}\bfseries Homework}\\[5pt]
        {\color{skymid}Complete the homework worksheet for Lesson 1.2:}\\[5pt]
        \begin{itemize}
          \setlength{\itemsep}{0.35em}\setlength{\topsep}{0em}
          \item {\color{white}NBA player dataset: classify all variables
            with justification}
          \item {\color{white}AP-style multiple-choice practice}
          \item {\color{white}Optional extension: real dataset from
            CDC or Census}
        \end{itemize}
        \vspace{0.45cm}
        {\color{skymid}\itshape\small
          \textbf{Preview:} Lesson 1.3 --- now that we know a variable is
          categorical, how do we display and summarize it?}
      \end{column}

    \end{columns}
  \end{minipage}
\end{frame}
}

\end{document}
